\section{Deforming representations: another way to parametrize the components}\label{deformation}

Having found a specific representation $\rho: \Gamma_8 \to \mathrm{SL(3,\mathbb C)}$ by
algebraic means, there is a complementary method we can often use to determine
the component $\mathcal{V}$ of the variety containing $\rho$. We apply this procedure to a representation in each of the three components found in the previous section. As a product, we get an explicit parametrisation of the representations, i.e. an explicit parametrisation of matrices generating the representation, at least for a Zariski-open subset of the studied component. This information is of course very valuable for the study of the representations found, as in \cite{DF}.


The procedure may be broken
down into four steps, which we now describe briefly; a detailed account of the
method is given in \cite{CLTb}.  Repeated use is made of the LLL algorithm \cite{LLL}
for finding an algebraic number of low degree, close to a number given numerically to
high precision.  For ease of exposition, we shall assume that
$\mathcal{V}$ has dimension $2$; in advance we might not know this dimension,
but usually with experimentation it quickly becomes evident.  Also, the dimension of
the Zariski tangent space at $\rho$ (\cite{CLTa}) provides an upper bound for the
dimension of $\mathcal{V}$.

\begin{itemize}
\item[\bf (i)]  First, we perturb the images of the group generators under $\rho$
slightly, and use these perturbed matrices as a starting point for Newton's method
in conjunction with the group relators, so as to converge numerically to a
representation $\sigma$.  If the original representation $\rho$ is not isolated,
$\sigma$ will typically be some generic representation close to $\rho$; however,
by carefully imposing extra constraints one can ``steer'' the Newton process so
that $\sigma$ has some specified character.  In this way we can obtain a
finite array of representations $\{\sigma_{ij}\}$ such that for each generator $g_k$ of $\Gamma_8$
the traces of the matrices $\alpha_{ij} = \sigma_{ij}(g_k) $ form a suitable array of (numerical
approximations to) rational numbers.  Of course the matrix $\alpha_{ij}$ depends on the
index $k$ of $g_k$, but to avoid excess of notation we omit mention of this index.
\item[\bf (ii)]  Matrices $\beta_{ij}$ are chosen (independently of $k$) so that the conjugates
$\alpha'_{ij} = \beta^{-1}_{ij} \alpha_{ij} \beta_{ij}$ have entries that are all algebraic.
This is always possible, owing to the fact that $\mathcal{V}$ is an algebraic
set; in practice we desire that the degrees of these algebraic numbers should be
as small as possible.
\item[\bf (iii)]  So far, for each group generator $g_k$ we have a two-dimensional array
of matrices $\alpha'_{ij}$ whose entries
are numerical approximations to elements of a number field $\mathcal{F}_{ij}$.
We choose a suitable basis for $\mathcal{F}_{ij}$ over the field of rationals,
noting that the field depends on the parameters $i, j$.  Next, LLL allows us to emerge
from the world of numerical approximations to that of exact quantities by guessing
an expression for each matrix entry as a linear combination of basis elements.
Of course we wish also to be freed from the discrete coordinates $(i, j)$ to
general (complex) coordinates $(u, v)$ for the variety $\mathcal V$, and this is
achieved, up to judicious guesswork, by means of polynomial interpolation across
the $(i,j)$--lattice.

\item[\bf (iv)]  The reader will note that guesswork has been employed at several
points of this procedure, and that substantial use was made of numerical approximations
to algebraic numbers.  However, we have arrived at a function $\Phi$ that associates to
each generator $g_k$ of $\Gamma_8$ an exact matrix, whose entries lie in an algebraic extension
of $\mathbb Q(u,v)$, where $u,v$ are independent transcendentals.  Using software
such as Mathematica or Maple, it is relatively straightforward to verify by formal
manipulation that the matrices $\Phi(g_k)$ satisfy the group relations, {\it i.e.} that
$\Phi$ extends to a representation of $\Gamma_8$.  Following the nomenclature of M. Culler and
P. Shalen, we call $\Phi$ the {\it tautological represention} of $\Gamma_8$ associated to $\mathcal{V}$.
A tautological representation depends of course on a chosen parametrization, {\it i.e.}
a coordinate system for the variety; also, two matrix representations are considered
to be equivalent if they agree up to postmultiplication by an inner automorphism of
the target linear group.
\end{itemize}

The method just described would appear at first sight to be remarkably {\it ad hoc;}\/
however, it often works well.  It was applied to three representations
$\rho_1 \,,\, \rho_2 \,,\, \rho_3$ given in \S \ref{section:example-unipotent} -- one in each of the components found in the previous section. It produces tautological
representations $R_i \; (i=1,2,3)$ of the figure-eight knot group
$\Gamma_8$ into $\mathrm{SL(3,\mathcal{F})}$, where the field $\mathcal{F}$ is an
algebraic extension of $\mathbb Q(u,v)$.  Our brief descriptions of the salient
properties of the $R_i$ are given mostly without proof; however, the reader can
check most of them with the aid of suitable computer algebra software.
\vvv\noind

\subsection{The first component, $R_1$}  \label{ss:R1}
We use the presentation
\[ \Gamma_8 \h=\h \langle a \bc b \;|\; a\,b\,b\,b\,a\,B\,A\,A\,B \rangle \h,\]
where uppercase denotes inverse (see section \ref{section:intro-eight}). The group $\Gamma_8$ is generated by meridians $a\,b \bc a\,b\,b$,
and we may take
\[ m = a\,b \h,\h \ell = a\,b\,A\,B\,a\,B\,A\,b \]
as commuting (meridian, longitude) pair.
Our choice of parameters for the variety is
\[ u = \tr(A) \quad,\quad v = \tr(b) \h,\]
this choice being justified by the fact that
all matrix entries for the tautological representation lie in $\mathbb Q(u,v)$ itself:


\[
a \;\mapsto\; \left[\begin{array}{ccc}0&0&1\\1&0&-u\\0&1&\frac{(-1+v)^2(1+v)}{u}\end{array}\right] \textrm{ and }\]
\[b\;\mapsto\; \left[\begin{array}{ccc}1&
\frac{-u^3(-2+v)+(-1+v)^3(1+v)}{u(-1+v)^2v}&
-\frac{u(-2+v)}{-1+v}\\
0&v&
\frac{-v+v^3}{u}\\
0&
\frac{u}{v-v^2}&
-1\end{array}\right] \quad. \]

For notational convenience, let us identify $a \bc b \bc A \bc B$ with their images
under $R_1$.  These elements have the following traces:
\[ \tr(a) = \frac{(-1+v)^2(1+v)}{u} \h,\h \tr(A) = u \h,\h tr(b) = v \h,\h tr(B) = v \h.\]
On can also compute that
\[ \tr(a\,b) \;=\; \frac{u^3 + (-1 + v)^3(1 + v)^2}{uv(-1 + v)^2} \h,\]
from which it is evident that attempting to assign certain values to the traces
of $A \bc b$ results in singularities.
\vv

The classical hyperbolic representations are those satisfying the
constraint $\tr(a) = \tr(A)$, and one can compute that there are two {\it boundary unipotent}
hyperbolic representations, {\it i.e.} representations mapping peripheral elements
to unipotent matrices.  These are as follows:
\[ (u , v) \;=\; \left(-3 + 2\sqrt{3}\,i \h,\h \frac{1}{2} + \frac{3}{2}\sqrt{3}\,i\right) \quad\textrm{and}\]
\[   (u , v) \;=\; \left(-3 - 2\sqrt{3}\,i \h,\h \frac{1}{2} - \frac{3}{2}\sqrt{3}\,i\right) \h.\]
These form a complex conjugate pair, corresponding to the complete hyperbolic structure
on the figure-eight knot complement.  The deformation space of the complete structure
has one complex dimension, as in the classical $\mathrm{PSL(2,\mathbb C)}$ case.


\subsection{The other two components, $R_2$ and $R_3$}\label{ss:R_2-R_3}

For the other two varieties, we choose the standard ``parabolic'' presentation
\[ \Gamma_8 \h=\h \langle a \bc b \;|\; A\,b\,a\,B\,a\,b\,A\,B\,a\,B \rangle \h,\]
and parameters
\[ u = \tr(A) = \tr(B) \quad,\quad v = \tr(a) = \tr(b) \h.\]
Some matrix entries
lie in a quadratic extension of $\mathbb Q(u,v)$, namely the field generated by
$\sqrt{\Delta}$, where
\[ \Delta \;=\; 4 u^3 + 4 v^3 - u^2 v^2 - 16 u v + 16 \h.\]
Either choice of square root of $\Delta$ yields a valid representation.
We shall see shortly that these varieties are related via precomposition with an
outer automorphism of the knot group.  Here are the tautological representations:
\vvv

\fbox{$R_2$}
\vskip-4ex

{\bfi{12}{21}\rm
\begin{align*}
   a \;\mapsto\; &\left[\begin{array}{ccc}
    \frac{v}{2} & 1 & -\frac{(1-i)(-16 + 8uv - 2v^3 - 4\sqrt{\Delta})}{8u^2 - 6uv^2 + v^4}\\
    \frac{1}{8}(1+i)(-2u+v^2) & \frac{1}{4}(1+i)v & 1\\
    \frac{1}{16}(8 - 4uv + v^3 - 2\sqrt{\Delta}) & \frac{1}{8}(-4u + v^2) & \frac{1}{4}(1-i)v
    \end{array}\right] \\[3ex]
   b\;\mapsto\; &\left[\begin{array}{ccc}
    \frac{v}{2} & i & \frac{(1+i)(-16 + 8uv - 2v^3 - 4\sqrt{\Delta})}{8u^2 - 6uv^2 + v^4}\\
    -\frac{1}{8}(1+i)(-2u+v^2) & \frac{1}{4}(1-i)v & i\\
    -\frac{1}{16}(8 - 4uv + v^3 - 2\sqrt{\Delta}) & -\frac{i}{8}(-4u + v^2) & \frac{1}{4}(1+i)v
    \end{array}\right]
\end{align*}
}

\fbox{$R_3$}
\vskip-4ex

{\bfi{12}{21}\rm
\begin{align*}
   a \;\mapsto\; &\left[\begin{array}{ccc}
    \frac{v}{2} & 1 & -\frac{(1-i)(-16 + 8uv - 2v^3 + 4\sqrt{\Delta})}{8u^2 - 6uv^2 + v^4}\\
    \frac{1}{8}(1+i)(-2u+v^2) & \frac{1}{4}(1+i)v & 1\\
    \frac{1}{16}(8 - 4uv + v^3 + 2\sqrt{\Delta}) & \frac{1}{8}(-4u + v^2) & \frac{1}{4}(1-i)v
    \end{array}\right] \\[3ex]
   b\;\mapsto\; &\left[\begin{array}{ccc}
    \frac{v}{2} & 1 & -\frac{(1-i)(-16 + 8uv - 2v^3 + 4\sqrt{\Delta})}{8u^2 - 6uv^2 + v^4}\\
    \frac{1}{8}(1-i)(-2u+v^2) & \frac{1}{4}(1-i)v & -i\\
    \frac{i}{16}(8 - 4uv + v^3 + 2\sqrt{\Delta}) & \frac{i}{8}(-4u + v^2) & \frac{1}{4}(1+i)v
    \end{array}\right]
\end{align*}
}\ind
\vv

Our first observation is that $R_2$ and $R_3$ satisfy the same relationships with
regard to automorphisms of the knot group $\Gamma_8$ as those set out in Proposition 9.2
of \cite{DF}.  In particular, if $\tau: \Gamma_8 \to \Gamma_8$ is the (non-inner) automorphism
determined by the assignments $a \mapsto a^{-1}\,b\,a \;,\; b \mapsto b$, then
$R_3$ agrees with $R_2 \circ \tau$ up to conjugation by a matrix in
$\mathrm{GL(3,\mathcal F)}$.
\vv

Secondly, the image of $R_2$ (also that of $R_3$) is a copy of the $(3,3,4)$--triangle group
\[ T(3,3,4) \;=\; \langle \alpha \bc \beta \bc \gamma \;|\;
   \alpha^3 = \beta^3 = \gamma^4 = \alpha\beta\gamma = 1 \rangle \h.\]
Specifically, if $c$ denotes the commutator $a\,b\,A\,B$\,, we have
\[ R_2(c) \;=\; \left[\begin{array}{ccc}1&0&0\\0&i&0\\0&0&-i\end{array}\right] \h,\]
independently of $u \bc v$, from which it follows immediately that $R_2(c)$ has order $4$.
Moreover, one can verify that the images under $R_2$ of $a\,c \;,\; c\,a\,c$ each have
order $3$.  Since $a \bc c$ generate the knot group $\Gamma_8$, the result follows by taking
$\alpha = R_2(a\,c) \;,\; \beta = R_2(c\,a\,c)^{-1} \;,\; \gamma = R_2(c)$.  It was
previously known that the images of the representations $\rho_2 \bc \rho_3$
were isomorphic to $T(3,3,4)$, see \cite{DF}.

\subsection{A focus on the $\mathrm{PU}(2,1)$ case}

As studied e.g. in \cite{DF}, the representations with image inside the real group $\mathrm{PU}(2,1)$ may carry interesting geometric properties. So we seek representations preserving Hermitian forms.

In the case of the first component $R_1$, the following condition on traces holds:
\[ \tr(A) \,=\, \overline{\tr(a)} \quad,\quad \tr(B) \,=\, \overline{\tr(b)} \h.\]
In particular, we note that since $\tr(B) = \tr(b) = v$, the parameter $v$ is forced to be real
for such representations.  Then, from $\tr(A) = \tr(a)$, we obtain
\[ 
\frac{(-1+v)^2(1+v)}{\overline{u}} \;=\; u\;,
\quad\text{or}\quad
u \;=\; |  \left(-1+v\right)\sqrt{1+v} |  \, z ,
\]
with $z$ lying on the unit circle.
\vv

It is in fact possible to compute a Hermitian matrix $H$ as a function of the coordinates
$v \bc z$,
\[ H(v,z) \;=\; \left[\begin{array}{ccc}1&r&s\\\overline{r}&1&r\\\overline{s}&\overline{r}&1\end{array}\right]\h,\]
whose associated form is respected by the representations just described;
the entries of $H$ lie in an Abelian extension of degree $4$ of $\mathbb Q(v,\mathrm{Re}(z))$,
and are a little too cumbersome for inclusion here.  However, for specific $v \bc z$
the signature of $H$ enables one to determine whether the representation
takes values in $\mathrm{SU(3)}$ or $\mathrm{SU(2,1)}$, or whether the associated
form is degenerate.  Indeed, the representations with degenerate forms occupy some
isolated points in the $(v,z)$--variety, together with some simple closed curves that
separate $\mathrm{SU(3)}$--representations from $\mathrm{SU(2,1)}$--representations.
\vv

One can compute that the variety $R_1$ gives rise to a single boundary
unipotent representation into $\mathrm{PU(2,1)}$, obtained by composing the representation
\[ (u , v) \;=\; (-\sqrt{3}\,i \;,\; 2) \]
with the natural projectivization map $\mathrm{SU(2,1) \to PU(2,1)}$.  The variety
$R_1$ also contains the other two lifts of this $\mathrm{PU(2,1)}$--
representation to $\mathrm{SU(2,1)}$, at
\[ (u , v) \;=\; (-\omega(\sqrt{3}\,i) \;,\; 2) \, ,\quad
   (u , v) \;=\; (-\omega^2(\sqrt{3}\,i) \;,\; 2) \quad
   (\text{where}\quad \omega = e^{2\pi i/3}).\]

\bigskip

The other two components $R_2$ and $R_3$ each give rise to just one boundary unipotent representation into
$\mathrm{PU(2,1)}$, given by $(u,v) = (3,3)$.
\vv

A necessary condition on traces for a representation in $R_2$ or $R_3$ to respect a Hermitian
forms is
\[ \tr(a) = \tr(A) \;,\; \tr(b) = \tr(B) \quad,\quad\text{\it i.e.}\quad v = \overline{u} \h.\]

Using the same technique as that described at the beginning of this section,
{\it i.e.} invoking LLL together with polynomial interpolation, the following diagonal matrix:
\[ H \;=\; \left[\begin{array}{ccc}
-\frac{1}{8}(\Delta - 16)\left(\sqrt{\Delta} + |u|^2 - 4 \right) & 0 & 0\\
0 & \Delta - 16 & 0\\
0 & 0 & -8\left(\sqrt{\Delta} + 4 \right)
\end{array}\right] \]
was found, satisfying
\[ H a^{-1} \;=\; a^{\,\dagger}H \h,\h H b^{-1} \;=\; b^{\,\dagger}H
\quad (v = \overline{u} \;,\; \Delta \geq 0) \h,\]
where the superscript $^\dagger$ denotes conjugate transpose, and
$\Delta = 4 u^3 + 4 v^3 - u^2 v^2 - 16 u v + 16 \;=\;
4(u^3 + \overline{u}^3) - |u|^4 - 16|u|^2 + 16$.
Clearly the condition $\Delta \geq 0$ is required for $H$ to be Hermitian.  We note that
the form represented by $H$ degenerates at $\Delta = 16$; in the present context $v = \overline{u}$
the corresponding values of $u$ are $0 \,,\, 4 \,,\, 4e^{2\pi i/3} \,,\, 4e^{4\pi i/3}$.
\vv

Substituting $u = x + iy \bc v = x - iy \h (x \bc y \in \mathbb R)$ in the expression for
$\Delta(u,v)$, we obtain (with mild abuse of notation)
\[ \Delta(x,y) \;=\; -x^4 - y^4 - 2x^2 y^2 - 24x y^2 + 8x^3 - 16x^2 - 16y^2 + 16 \h.\]
\ind

We have therefore established that
representations respecting Hermitian forms occur in the bounded complementary regions
in the $(x,y)$--plane of the curve $\Delta(x,y) = 0$, see Figure \ref{fig:su(2,1)}.
The $\mathrm{SU(3)}$ representations lie in the central complementary region, and
the $\mathrm{SU(2,1)}$ representations lie in the three petals.  Corresponding
Hermitian representations occur also on the curve itself, except for singularities
at the three self-intersection points of the curve, where representations are
not defined.  The quantity $\Delta$ has
geometric significance in that it is related to the Brehm shape invariant \cite{Brehm}
of the triangle in the associated homogeneous space (complex projective plane
$\mathbb{CP}^2$ or complex hyperbolic plane $\mathbb{CH}^2$), whose vertices are the
fixed points of the elliptic isometries $\alpha \bc \beta \bc \gamma$.

\begin{figure}[ht]
\includegraphics[scale=1]{r2form-eps-converted-to.pdf}
\caption{$\mathrm{SU(3)}$ and $\mathrm{SU(2,1)}$ representations in $R_2$}\label{fig:su(2,1)}
\end{figure}



\bibliography{bibli-1}
\bibliographystyle{plain}




\end{document}

\subsection{ANY DIMENSION ???}


\begin{prop}
Let N be a compact orientable 3-manifold. Let $\rho_0 : \pi_1(N)\rightarrow \SL(n,\bC)$ be an irreducible representation such that for each torus component  of $\partial N$, the image of ${\rho_{_0}}_{ \vert\pi_1(\partial N)}$ contains a unipotent element with only one eigenvector.  Let
$X_0$ be an irreducible component of $X (N)$ containing $[\rho_0]$.  
Then ${\rm dim} X_0 \geq (n-1)s - (n^2-1)\chi (N)$, where $s$ is the number of torus components in $\partial N$.
\end{prop}

\begin{pf}
We follow Culler-Shalen.  Working in the representation variety, we need to prove that ${\rm dim} H_0 \geq (n-1)s - (n^2-1)\chi (N)+(n^2-1)$
where $H_0\subset Hom(\pi_1(N),SL(n,\bC)$ is the irreducible component containing $\rho_0$. 

The proof is by induction on $s$.  Suppose first that $s=0$.  We realize $N$ as a finite 2-dimensional CW-complex with one 0-cell, $m$ 1-cells and 
$n$ 2-cells so that $\pi_1 =\langle \ g_1, \cdots , g_m \ \vert\ r_1\cdots r_n\ \rangle$ is a presentation of the fundamental group. Define the regular map $f:\SL(n,\bC)^m\rightarrow \SL(n,\bC)^n$ via the r relators.  The representation variety is the inverse image $f^{-1}(1,\cdots ,1)$ and therefore
any irreducible component has dimension  
$$
\geq  (n^2-1)m-(n^2-1)n=- (n^2-1)\chi (N)+(n^2-1).
$$
 Suppose now $s> 0$ and let $T$ be a torus component.  Let $\alpha\in \pi_1(T)$ be an element  such that  $\rho_0(\alpha)$ is unipotent.

 
 Let $\gamma\in \pi_1(N)$, with base point in $T$, be the corresponding element as in the lemma.   We can drill a torus following a tubular neighborhood along $\gamma$ and obtain as in Culler-Shalen a manifold $N'$ where the boundary torus is transformed into a genus 2 surface.
 The standard basis $\alpha',\beta',\gamma',\delta'$ is such that $\alpha'$ and $\gamma'$ are mapped to $\alpha$ and $\gamma$
 respectively and $\delta'$ is the meridian of the drilled torus whose normal closure is the kernel of the embedding $i_* :\pi_1(N')\rightarrow  \pi_1(N)$.
 
 We apply the induction hypothesis to the irreducible component $H'_0$ of containing $\rho_0\circ i_*\in  Hom(\pi_1(N),SL(n,\bC))$ to obtain
 
 $$
dim H'_0 \geq (n-1)(s-1) - (n^2-1)\chi (N')+ (n^2-1).
$$
$$
= (n-1)s - (n^2-1)\chi (N)+ (n^2-1)+n(n-1).
$$
Observe that $Hom(\pi_1(N),SL(n,\bC)$ is identified to the closed variety  $W\subset Hom(\pi_1(N'),SL(n,\bC)$ defined by the representations
with $\rho(\delta')=1$.

We will show that $H_0$ is an irreducible component of the inverse image of a point under a regular map $f: H'_0\rightarrow \bC^{n(n-1)}$ and this will prove the proposition.  
\end{pf}
